\section{Appendix C: Translation bounds and hidden-allocation contraction}

The conditional response law in \cref{obj:lem:block-law} connects the complete observation experiment to translations of the block baseline density. The bounds below quantify the information carried by these translations and by the remaining source allocation. Throughout, the retained graph retains its labels and detailed edge subsets, and averaging uses its actual marginal law.

\subsection{Baseline translations and common-marginal comparisons}

The first ingredient measures how a change in the block response location affects its square-root density. We fix the density and derivative conventions before stating the translation bound.



The square-root density in \cref{obj:lem:baseline-translation-affinity} supplies the geometry for comparing translated block responses. To carry these comparisons through conditioning and mixing, we also fix the probability-distance conventions.



These conventions allow one bound to address a single translated response, a vector of independent translated responses, and conditional response laws sharing the same marginal for the conditioning variables. The following statement collects the corresponding comparisons.

\begin{lemmav}{lem:baseline-translation-affinity}[Baseline translation and mixture bounds]
Let \(f\) be the baseline density and \(\varphi\) its square root:
\[
f(w)=
\begin{cases}
4\cos^2(2\pi w),& |w|\le 1/4,\\
0,& |w|>1/4,
\end{cases}
\qquad
\varphi(w)=\sqrt{f(w)}.
\]
Then \(\varphi\) is globally \(4\pi\)-Lipschitz, and
\[
\int_{-1/4}^{1/4}\bigl(4\pi\sin(2\pi w)\bigr)^2\,dw
=4\pi^2.
\]
For every \(s,t\in\mathbb R\),
\[
\int_{\mathbb R}
\bigl(\sqrt{f(w-s)}-\sqrt{f(w-t)}\bigr)^2\,dw
\le 4\pi^2(s-t)^2.
\]
For every nonnegative integer \(B\) and every pair of vectors
\(s,t\in\mathbb R^B\),
\[
\int_{\mathbb R^B}
\left(
\sqrt{\prod_{\ell=1}^{B}f(y_\ell-s_\ell)}
-
\sqrt{\prod_{\ell=1}^{B}f(y_\ell-t_\ell)}
\right)^2\,dy
\le
\sum_{\ell=1}^{B}4\pi^2(s_\ell-t_\ell)^2.
\]
Here real and Euclidean integrals use Lebesgue measure, including the usual empty-product convention when \(B=0\).

For probability laws, let \(\operatorname{TV}\) denote total variation, defined as the supremum of the absolute difference of probabilities over measurable events. The unhalved squared Hellinger distance between laws with densities \(v_+\) and \(v_-\) relative to a common measure \(\mu\) is
\[
\int_\Omega
\bigl(\sqrt{v_+(x)}-\sqrt{v_-(x)}\bigr)^2\,d\mu(x).
\]
For every measurable space \(\Omega\), every sigma-finite measure \(\mu\), and every pair of nonnegative measurable densities \(v_+,v_-\) satisfying
\[
\int_\Omega v_+(x)\,d\mu(x)
=
\int_\Omega v_-(x)\,d\mu(x)
=1,
\]
one has
\[
\operatorname{TV}\bigl(v_+\,\mu,v_-\,\mu\bigr)
\le
\left[
\int_\Omega
\bigl(\sqrt{v_+(x)}-\sqrt{v_-(x)}\bigr)^2\,d\mu(x)
\right]^{1/2}.
\]

For the common-marginal version, let \(\Xi,\Omega\) be measurable spaces and suppose:
\begin{itemize}
\item \textbf{(Reference measures.)} \(\nu\) is a probability measure on \(\Xi\), and \(\mu\) is a sigma-finite measure on \(\Omega\).
\item \textbf{(Conditional densities.)} \(v_+,v_-:\Xi\times\Omega\to\mathbb R\) are jointly measurable and nonnegative.
\item \textbf{(Normalization.)} For every \(\xi\in\Xi\),
\[
\int_\Omega v_+(\xi,x)\,d\mu(x)
=
\int_\Omega v_-(\xi,x)\,d\mu(x)
=1.
\]
\end{itemize}
The joint laws are
\[
\nu(d\xi)\,v_\pm(\xi,x)\,\mu(dx),
\]
and their response mixtures are
\[
\int_\Xi v_\pm(\xi,x)\,\mu(dx)\,\nu(d\xi).
\]
The unhalved squared Hellinger distance between these joint laws equals
\[
\int_\Xi
\left[
\int_\Omega
\bigl(\sqrt{v_+(\xi,x)}-\sqrt{v_-(\xi,x)}\bigr)^2\,d\mu(x)
\right]d\nu(\xi),
\]
and the response mixtures satisfy
\[
\operatorname{TV}\left(
\int_\Xi v_+(\xi,x)\,\mu(dx)\,\nu(d\xi),
\int_\Xi v_-(\xi,x)\,\mu(dx)\,\nu(d\xi)
\right)
\le
\left[
\int_\Xi
\left[
\int_\Omega
\bigl(\sqrt{v_+(\xi,x)}-\sqrt{v_-(\xi,x)}\bigr)^2\,d\mu(x)
\right]d\nu(\xi)
\right]^{1/2}.
\]

Define the derivative representative, with value zero at the two support endpoints, by
\[
\varphi'(w)=
\begin{cases}
-4\pi\sin(2\pi w),& |w|<1/4,\\
0,& |w|\ge 1/4.
\end{cases}
\]
At every \(w\in\mathbb R\setminus\{-1/4,1/4\}\), \(\varphi\) is differentiable with derivative \(\varphi'(w)\), and
\[
\int_{\mathbb R}\varphi'(w)^2\,dw=4\pi^2.
\]
\end{lemmav}

The translation inequality in \cref{obj:lem:baseline-translation-affinity} gives a quadratic bound in the change of location. Its product version adds the squared changes across distinct block responses. The common-marginal identity then preserves the conditioning variables in the joint experiment: conditional squared Hellinger distances are averaged under their shared marginal. The response-mixture inequality supplies the corresponding comparison after those variables are integrated out. These are the distance comparisons used to translate conditional response bounds into testing bounds.

\subsection{The translated channel and all-order Walsh energy}

For a block with \(d\) source signs, the response translation depends on their sum. A uniform sign mixture provides a reference density against which every translated response can be expanded. We first define this family at a fixed admissible amplitude.



The normalized translated densities in \cref{obj:lem:walsh-channel-energy} admit an expansion in products of source signs. Symmetry makes the coefficient depend on the number of signs in the product. The same result records the coefficients and the two energy aggregates needed for the allocation bound.



In \cref{obj:lem:walsh-channel-energy}, the binomial weight counts the sign subsets of each order. The additional factor \(s\) in the order-weighted aggregate records their sizes. The result controls both aggregates while retaining every nonconstant order of the response expansion.

\begin{lemmav}{lem:walsh-channel-energy}[Walsh channel energy bound]
Let \(d\ge1\) be an integer and let \(0\le h\le1/4\). With \(f\) denoting the baseline density, define, for each sign vector \(\varepsilon\in\{-1,1\}^d\),
\[
F_\varepsilon(w)
=f\!\left(w-\frac{h}{2d}\sum_{j=1}^d\varepsilon_j\right),
\qquad
g_d(w)=2^{-d}\sum_{\varepsilon\in\{-1,1\}^d}F_\varepsilon(w).
\]
For every integer \(s\ge0\), define the Walsh coefficient function by
\[
a_s(w)=
\begin{cases}
2^{-d}\displaystyle\sum_{\varepsilon\in\{-1,1\}^d}
\frac{F_\varepsilon(w)}{g_d(w)}
\displaystyle\prod_{j=1}^{\min\{s,d\}}\varepsilon_j,
& g_d(w)\ne0,\\
0,&g_d(w)=0.
\end{cases}
\]
The empty product equals one. Define the integrated Walsh energies and their two aggregates by
\[
\gamma_s=\int_{\mathbb R}a_s(w)^2g_d(w)\,dw
\quad(1\le s\le d),
\qquad
\eta=\sum_{s=1}^d\binom ds\gamma_s,
\qquad
\eta_1=\sum_{s=1}^d s\binom ds\gamma_s.
\]
For every \(w\) with \(g_d(w)>0\), \(a_0(w)=1\), and for every \(\varepsilon\in\{-1,1\}^d\),
\[
\frac{F_\varepsilon(w)}{g_d(w)}
=
\sum_{E\subseteq\{1,\ldots,d\}}
a_{|E|}(w)\prod_{j\in E}\varepsilon_j.
\]
Moreover, for every integer \(s\ge1\),
\[
\int_{\mathbb R}a_s(w)g_d(w)\,dw=0,
\]
and the aggregate energies satisfy
\[
0\le\eta\le\eta_1\le\frac{12\pi^2h^2}{d}.
\]
\end{lemmav}

\Cref{obj:lem:walsh-channel-energy} supplies an exact expansion of the normalized response density on the positive-density region of the mixture. Its constant coefficient is one, and each nonconstant coefficient has mean zero under that mixture. The aggregate bound quantifies the squared size of all nonconstant terms, including their order weights, at scale \(h^2/d\). This gives a single measure of response information that can be used across the possible hidden-source allocations.

\subsection{Conditional reference laws and hidden allocation}

The complete-record comparison combines response information with the source associations revealed by the audit. Conditional on the retained graph, the reference law preserves the actual assignment marginal and replaces the distinct block responses by independent uniform sign-mixture responses. The required capacities and conditional discrepancies are defined together.



The reference law in \cref{obj:lem:hidden-allocation-contraction} keeps every assignment coordinate in the comparison. The reconstruction property in \cref{obj:lem:block-law} connects the distinct responses to the full outcome record, including their repeated recipient coordinates. The remaining-capacity event measures how much source allocation is still hidden after observing the graph. Under a smallness condition on the total Walsh energy, that contraction result bounds the conditional discrepancy averaged over the complete graph.

\begin{lemmav}{lem:hidden-allocation-contraction}[Hidden allocation contraction bound]
In the complete block experiment, let \(n\) be the population size, \(B\) the number of blocks, \(d\) the number of sources per block, and \(m=Bd\). Suppose:
\begin{itemize}
\item \textbf{(Block sizes.)} \(n,B,d\) are integers satisfying \(n\ge4\), \(B\ge1\), \(d\ge1\), and \(2Bd\le n\).
\item \textbf{(Amplitude and audit.)} The amplitude satisfies \(0\le h\le1/4\). The retention probability satisfies \(0\le q\le1\); the audit marks on all ordered off-diagonal pairs are independent Bernoulli variables with success probability \(q\), independent of the assignment.
\item \textbf{(Assignment.)} The assignment law satisfies \cref{obj:ass:assignment}.
\item \textbf{(Channel smallness.)} Let \(\gamma_s\) be the integrated squared Walsh coefficient of order \(s\) for the block-response channel at amplitude \(h\), and define its total and order-weighted nonconstant energies by
\[
\eta=\sum_{s=1}^{d}\binom{d}{s}\gamma_s,
\qquad
\eta_1=\sum_{s=1}^{d}s\binom{d}{s}\gamma_s.
\]
These satisfy \(4e\eta<1\).
\end{itemize}
Let \(H\) be the complete labeled retained graph. For each block, let \(u_\ell\) be its remaining hidden-source capacity, and define the total remaining capacity and source-association revelation probability by
\[
M=\sum_{\ell=1}^{B}u_\ell,
\qquad
p=1-(1-q)^d.
\]
Conditional on \(H\), let \(Q_H\) have the actual assignment marginal and independent distinct block responses with the uniform sign-mixture response density \(g_d\). Let \(L_+\) be the density ratio of the positive-prior conditional law \(P_{+1,h}(Z,y\mid H)\) relative to \(Q_H\), and define
\[
\chi_H^2=\mathbb E_{Q_H}\!\left[(L_+-1)^2\right].
\]
Then
\[
\mathbb E_H\!\left[\mathbf1\{M\ge m/4\}\chi_H^2\right]
\le
\frac{\exp(eBp\eta_1)}{1-4e\eta}-1.
\]
The expectation is taken under the actual marginal law of the complete retained graph, including its labels and detailed retained-edge subsets.
\end{lemmav}

The block-size conditions in \cref{obj:lem:hidden-allocation-contraction} place the source and recipient groups within the finite population. Independent auditing and the assignment condition preserve the experimental design used by the conditional response law. The channel-smallness condition keeps the denominator positive and specifies the range in which the contraction bound applies.

The bound combines two aspects of the experiment. The exponential term depends on the revelation probability and the order-weighted response energy, while the denominator depends on the total nonconstant energy. The indicator restricts the averaged discrepancy to graph realizations leaving at least one quarter of the original source capacity hidden. All graph labels and detailed retained subsets remain part of the averaging law, so the comparison applies to the information represented by the complete retained graph.

Together, \cref{obj:lem:baseline-translation-affinity,obj:lem:walsh-channel-energy,obj:lem:hidden-allocation-contraction} provide the distance comparisons, response-energy control, and allocation bound used in \cref{obj:thm:block-testing-scale,obj:thm:nonlocal-testing-answer}. Their roles are complementary: translations control changes in block responses, Walsh energies quantify their dependence on source signs, and the conditional-reference comparison accounts for the hidden allocation under the actual observation design.
